
\documentclass[10pt]{article}
\usepackage[utf8]{inputenc}
\usepackage[T1]{fontenc}
\usepackage{amsmath, amssymb, xcolor, geometry, enumitem, multicol, tcolorbox}

\definecolor{mainblue}{RGB}{0, 102, 204}
\geometry{a4paper, margin=1.2cm, top=1cm, bottom=3cm, footskip=1.2cm}

\begin{document}
\enlargethispage{2.5cm}

% Modern Header
\begin{tcolorbox}[colback=mainblue!5, colframe=mainblue, sharp corners, boxrule=0.5pt]
    \small \textbf{Unit:} undefined \hfill \textbf{Name:} \rule{3cm}{0.4pt} \\
    \small \textbf{Topic:} Variables and Expressions / Order of Operations \hfill \textbf{Date:} \rule{2cm}{0.4pt} \\
    \centering \textbf{\large Foundation (Level -1)}
\end{tcolorbox}

\vspace{0.2cm}

% MCQ Section
\begin{multicols}{2}
\begin{enumerate}[label=\textbf{Q\arabic*.}, leftmargin=*, itemsep=6pt, parsep=0pt]

  \item \small What is the variable in the expression $2x$?
  \begin{enumerate}[label=(\Alph*), leftmargin=0.6cm, nosep, itemsep=2pt]
    \item 2
    \item x
    \item 2x
    \item +
    \item $-$
  \end{enumerate}

  \item \small Identify the coefficient in the expression $4y$.
  \begin{enumerate}[label=(\Alph*), leftmargin=0.6cm, nosep, itemsep=2pt]
    \item y
    \item 4
    \item 4y
    \item $+$
    \item $-$
  \end{enumerate}

  \item \small What is the operator in the expression $3 + 2$?
  \begin{enumerate}[label=(\Alph*), leftmargin=0.6cm, nosep, itemsep=2pt]
    \item 3
    \item 2
    \item $+$
    \item $-$
    \item $	imes$
  \end{enumerate}

  \item \small Identify the constant term in the expression $2x + 5$.
  \begin{enumerate}[label=(\Alph*), leftmargin=0.6cm, nosep, itemsep=2pt]
    \item 2x
    \item 5
    \item $+$
    \item $-$
    \item $	imes$
  \end{enumerate}

  \item \small What is the expression $x$ an example of?
  \begin{enumerate}[label=(\Alph*), leftmargin=0.6cm, nosep, itemsep=2pt]
    \item A constant
    \item A variable
    \item A coefficient
    \item An operator
    \item A constant term
  \end{enumerate}

  \item \small Identify the type of expression $3x$ is.
  \begin{enumerate}[label=(\Alph*), leftmargin=0.6cm, nosep, itemsep=2pt]
    \item A numerical expression
    \item An algebraic expression
    \item A constant expression
    \item A variable expression
    \item A geometric expression
  \end{enumerate}

  \item \small What does the term $2x$ represent in an expression?
  \begin{enumerate}[label=(\Alph*), leftmargin=0.6cm, nosep, itemsep=2pt]
    \item The sum of 2 and x
    \item The product of 2 and x
    \item The difference of 2 and x
    \item The quotient of 2 and x
    \item The constant term
  \end{enumerate}

  \item \small Which of the following is an example of a simple expression?
  \begin{enumerate}[label=(\Alph*), leftmargin=0.6cm, nosep, itemsep=2pt]
    \item $2x + 3$
    \item $4y - 2$
    \item $x$
    \item $5$
    \item All of the above
  \end{enumerate}

\end{enumerate}
\end{multicols}

\vspace{-0.2cm}
\hrule
\vspace{0.2cm}
\textbf{\small Open Ended Questions (FRQ)}
\begin{enumerate}[label=\textbf{Q\arabic*.}, start=9, itemsep=5pt, topsep=0pt]

  \item \small Draw a simple diagram to illustrate the concept of a variable in an expression. Explain your diagram in 1-2 sentences. \vspace{1.2000000000000002cm}

  \item \small Provide an example of a real-life situation where you would use a simple expression, such as $2x$ or $3y$. Explain your example in 1-2 sentences. \vspace{1.2000000000000002cm}

\end{enumerate}

\vfill

% Chef's Tips Footer
\begin{tcolorbox}[colback=blue!5, colframe=mainblue!50, title=\small \textbf{Chef's Tips}, fonttitle=\bfseries, bottom=1mm]
\begin{itemize}[noitemsep, topsep=2pt, leftmargin=0.5cm]
  \item \scriptsize Emphasize the definition of variables as letters or symbols that represent unknown or changing values.\item \scriptsize Focus on simple expressions, such as $2x$ or $3y$, to introduce the concept of coefficients and variables.\item \scriptsize Use visual aids and real-life examples to help students understand the concept of expressions and variables.
\end{itemize}
\end{tcolorbox}

\end{document}
  